\documentclass[12pt,a4paper]{article}
\usepackage[spanish]{babel} % Corta palabras en espa~nol
\usepackage[utf8]{inputenc} % Escribir con acentos, ~n...
\usepackage{eurosym} % smbolo del euro


\begin{document}
%\secttion{Introducción}
\leftline{\textbf{\huge{Ejemplos de Listas y Tablas}}}
\medskip
Este primer ejemplo trata de demostrar la facilidad de
\LaTeX{} para hacer \textbf{listas}:

Ejemplo 1
\begin{itemize}
\item primer punto
\item segundo punto
\end{itemize}

Ejemplo 2
\begin{itemize}
\item[*] punto uno
\item[$*$] punto dos
\item[$\circ$] punto tres
\end{itemize}

Ejemplo 3
\begin{enumerate}
\item punto uno
\begin{enumerate}
\item pto uno de 1
\item pto dos de 1
\end{enumerate}
\item punto dos
\end{enumerate}

\textbf{Veamos ahora las tablas:}


\begin{tabular}{lrc}
Nombre & Edad & Clase \\
\hline
José & 24 & P \\
Juanito & 9 & P+ \\
Carlos & 11 & Q-
\end{tabular}

\textbf{Otra tabla}

\begin{center}
\begin{tabular}{|l||r|p{2cm}|}
\hline
Nombre & Edad & Clase \\
\hline \hline
José & 24 & El otro día
estaba en clase. \\
Juanito & 9 & P+ \\
\hline
\end{tabular}
\end{center}



\mbox{\textbf{En LATEX es posible generar cajas de varios tipos.}}

\fbox{Hola que tal}

fin = f\null{}in = f\mbox{}in?

\fbox{Esta \raisebox{-0.1cm}{forma}
\raisebox{-0.3cm}{de} \raisebox{-0.4cm}{escribir}
\raisebox{-0.6cm}{me} \raisebox{-0.8cm}{marea} un
\raisebox{0.1cm}{poco}}.


\end{document}